% from ega2-1-fr.tex
\oldpage[II]{188}
cohérent $\mathcal{I}$ d'idéaux de $\mathcal{O}_C$. Comme $p$ est séparé et
quasi-compact, $X$ est alors un schéma quasi-compact, et on est ramené à
prouver que $\mathcal{L}$ est ample
(\hyperref[II.4.6.4-fr]{4.6.4}). En vertu du critère
(\hyperref[II.4.5.2-fr]{4.5.2, a)}), on doit donc prouver qui ce suit~: pour
tout Idéal quasi-cohérent $\mathcal{K}$ de $\mathcal{O}_X$ et tout point
$x\in X$ n'appartenant pas au support de $\mathcal{O}_X/\mathcal{K}$, il
existe un entier $n>0$ et une section $f$ de
$\mathcal{K}\otimes\mathcal{L}^{\otimes n}$ au-dessus de $X$ telle que
$f(x)\neq0$.

Pour cela, posons
\[
  \mathcal{K}_V=\pi^*(\mathcal{K})\mathcal{O}_V
\]
et
\[
  \mathcal{K}_W=\mathcal{K}_V|W,\qquad\hbox{où}\qquad W=V-j(X)~;
\]
comme la restriction de $q$ à $W$ est une immersion quasi-compacte dans $C$,
il résulte de (\hyperref[I.9.4.2-fr]{I, 9.4.2}) que $\mathcal{K}_W$ est la
restriction à $W$ d'un Idéal quasi-cohérent $\mathcal{K}'_V$ de
$\mathcal{O}_V$, de la forme
\[
  \mathcal{K}'_V=q^*(\mathcal{K}_C)\mathcal{O}_V
\]
où $\mathcal{K}_C$ est un Idéal quasi-cohérent de $\mathcal{O}_C$. En outre,
comme par hypothèse $q^{-1}(Y)\subset j(X)$ et que $Y$ est défini par l'Idéal
$\mathcal{I}$, la restriction à $W$ de $q^*(\mathcal{I})\mathcal{O}_V$ est
identique à celle de $\mathcal{O}_V$, donc $\mathcal{K}_W$ est aussi la
restriction à $W$ de $q^*(\mathcal{I}\mathcal{K}_C)\mathcal{O}_V$, et l'on
peut donc supposer que l'on a $\mathcal{K}_C\subset\mathcal{I}$, d'où
\begin{equation}
\label{II.8.10.6.1-fr}
  \mathcal{K}'_V\subset q^*(\mathcal{I})\mathcal{O}_V\subset\mathcal{J}
\tag{8.10.6.1}
\end{equation}
compte tenu de (\hyperref[I.4.4.6-fr]{I, 4.4.6}) et de la commutativité de
(\hyperref[II.8.9.4.1-fr]{8.9.4.1}). En outre, on déduit de
(\hyperref[II.8.10.4-fr]{8.10.4}) que l'on a
\begin{equation}
\label{II.8.10.6.2-fr}
  \mathcal{K}'_V\subset\mathcal{K}_V.
\tag{8.10.6.2}
\end{equation}

Cela étant, il résulte de (\hyperref[II.8.10.3-fr]{8.10.3}) que $j(x)$
appartient à une seule composante irréductible de $\pi^{-1}(x)$~; soit $z$ le
point générique de cette composante, et posons $z'=q(z)$. En vertu de
(\hyperref[II.8.10.5-fr]{8.10.5}), la démonstration sera achevée (compte tenu
de (\hyperref[II.8.10.6.1-fr]{8.10.6.1}) et
(\hyperref[II.8.10.6.2-fr]{8.10.6.2})) si nous prouvons l'existence d'une
section $g$ de $\mathcal{K}'_V$ au-dessus de $V$ telle que $g(z)\neq0$. Or,
par hypothèse, $\mathcal{K}$ a une restriction égale à celle de
$\mathcal{O}_X$ dans un voisinage ouvert de $x$~; d'autre part, il résulte de
(\hyperref[II.8.10.3-fr]{8.10.3}) que l'on a $z\neq j(x)$, donc $z\in W$, et
par suite $(\mathcal{K}_W)_z=\mathcal{O}_{V,z}$, d'où par définition
$(\mathcal{K}_C)_{z'}=\mathcal{O}_{C,z'}$. Puisque $C$ est affine, il y a donc
une section $g'$ de $\mathcal{K}_C$ au-dessus de $C$ telle que
$g'(z')\neq0$, et en prenant pour $g$ la section de $\mathcal{K}'_V$
correspondant canoniquement à $g'$, on a bien $g(z)\neq0$, ce qui achève la
démonstration.
\end{env}

\begin{remark}[8.10.7]
\label{II.8.10.7-fr}
Nous ignorons si, dans l'énoncé de (\hyperref[II.8.9.4-fr]{8.9.4}), la
condition (ii) est ou non superflue. En tout cas, la conclusion ne subsiste
plus si l'on ne suppose pas l'existence d'un $Y$-morphisme $\pi:V\to X$ tel
que $\pi\circ j=1_X$~; indiquons brièvement comment on peut en effet former un
contre-exemple, dont les détails ne pourront être développés que plus tard.
On prend $Y=\operatorname{Spec}(k)$ où $k$ est un corps,
$C=\operatorname{Spec}(A)$, où $A=k[T_1,T_2]$, la $Y$-section $\varepsilon$
correspondant à l'homomorphisme d'augmentation $A\to k$. On désigne par $C'$
le schéma déduit de $C$ par éclatement du point fermé $a=\varepsilon(Y)$ de
$C$~; si $D$ est l'image réciproque de $a$ dans $C'$, on considère dans $D$
un point fermé $b$,

